% =====================================================================
\lezione{06/11/2026}{Laboratorio 4: conformity audit CRA--GDPR (Sestante Mobility)}
  {Gianluca Gilardi}
  {lavoro di gruppo sul caso Sestante Mobility; audit esterno di un fabbricante;
   registro dei finding, opinione e management letter; consegna tramite workbench online}
\label{lez:2026-11-06}
% =====================================================================

\begin{nota}[Fonte]
Trascrizione automatica della registrazione (lezione~14): la sessione è quasi interamente silenzio,
perché i gruppi lavorano in autonomia. Si sentono solo l'apertura, due avvisi sul tempo residuo e i
saluti finali. Il contenuto del caso, la consegna e i numeri provengono dal materiale del laboratorio:
\emph{Conformity\_Audit\_Playbook} e slide \emph{Conformity\_Audit\_Exercise\_Briefing}. Lo
svolgimento completo è nel Capitolo~\ref{lab:04}.
\end{nota}

\section{Sintesi}
Quarto laboratorio, presentato da Gianluca al termine della lezione del 02/11
(Capitolo~\ref{lez:2026-11-02}). Gli studenti lavorano da \textbf{auditor esterni} nominati dal
comitato di audit di \textbf{Sestante Mobility S.p.A.}, una media impresa che produce un box
telematico per flotte aziendali. Partendo da una \emph{evidence room} di 16 documenti e da due
interviste, devono stabilire se il fascicolo tecnico consente di tenere il prodotto sul mercato ai
sensi del CRA e quale ruolo ha l'azienda ai fini del GDPR. Il lavoro si chiude con un registro dei
finding, un'opinione di audit e una management letter. Rispetto alla simulazione Pulse K cambia la
prospettiva: non si gestisce un incidente dall'interno, si valuta a posteriori ciò che l'azienda può
dichiarare di sé.

\section{Consegna}
\begin{itemize}
  \item Si lavora sul \textbf{Conformity Audit Workbench}, uno strumento online che ricalca i cinque
    passi del playbook (scope, evidence room, registro, opinione, testo finale). Il deliverable è il
    testo prodotto al passo~5, copiato con l'apposito pulsante in un file di testo semplice, con il
    nome del gruppo e dei membri in cima. Un file per gruppo, caricato entro la fine della sessione.
  \item Il file ha tre blocchi: \textbf{A}~risposte di scope (quali obblighi CRA per quali unità;
    ruolo GDPR per tre flussi di dati), \textbf{B}~registro di almeno quattro finding di livello
    Medium o superiore, con un controllo che funziona e almeno un criterio esaminato e scartato,
    \textbf{C}~opinione e management letter. Bonus: mappatura di ogni finding su un controllo
    dell'Annex~A.8 della ISO/IEC~27001.
\end{itemize}

\section{Indicazioni udibili in aula}
\begin{itemize}
  \item Gianluca dà per scontato che il playbook sia stato studiato prima del laboratorio; il link
    allo strumento è stato condiviso in chat.
  \item \textbf{Nessuno deve completare l'audit}: l'obiettivo è arrivare ad \textbf{almeno quattro
    finding}, che per il docente sono ``più che sufficienti''.
  \item Il risultato può essere giusto o sbagliato e tutto è discutibile; ciò che non si discute è la
    \textbf{logica}. Serve un percorso logico coerente, come già suggerito dopo il laboratorio
    precedente.
  \item Le domande si fanno in chat; il microfono viene aperto su richiesta.
  \item Chiusura prevista verso le 13:15, come nel laboratorio precedente.
  \item La discussione degli elaborati era annunciata per la lezione del lunedì successivo. Il 09/11
    però nessun gruppo ha presentato e le presentazioni sono state rinviate
    (Capitolo~\ref{lez:2026-11-09}).
\end{itemize}

\begin{esame}
Il laboratorio applica le regole di metodo annunciate il 02/11: le date di applicazione del CRA
valgono \textbf{per singola unità} e non per modello; \textbf{una fonte è un'asserzione, due fonti
concordanti sono un fatto}; un certificato ISO/IEC~27001 non prova la conformità né al CRA né al GDPR;
ciò che non è nelle evidenze non si presume. Sono competenze direttamente utili per le domande
d'esame sulla qualificazione di un caso.
\end{esame}

\section{Domande aperte e punti da verificare}
\begin{verifica}
\begin{itemize}
  \item L'orario di chiusura (``1.15'') e i due avvisi sul tempo residuo sono ricavati dalla
    trascrizione; le slide indicano per il laboratorio la fascia 11:30--13:00.
  \item Il debrief in aula (cinque domande indicate nelle slide di briefing) non risulta ancora
    svolto: verificare se e quando verrà ripreso.
  \item Cognome del co-docente Gianluca: l'indirizzo di consegna riportato nel materiale del
    laboratorio fa pensare a ``Gilardi'', ma non è confermato in aula.
\end{itemize}
\end{verifica}
